\documentclass[11pt]{article}
\usepackage[a4paper,margin=25mm]{geometry}
\usepackage{amsmath,amssymb,amsthm,booktabs,hyperref,xurl,microtype}
\setlength{\emergencystretch}{3em}
\hypersetup{hidelinks}
\newtheorem{theorem}{Theorem}
\newtheorem{lemma}[theorem]{Lemma}
\newtheorem{corollary}[theorem]{Corollary}
\newcommand{\Q}{\mathbb Q}\newcommand{\Z}{\mathbb Z}
\title{Fixed dyadic translations for simultaneous contact certificates}
\author{\(\mathbb Z\mathbb Q\) Diophantine Project}
\date{23 November 2026}
\begin{document}
\maketitle
\begin{abstract}
The ordinary-abc quartic interface has one contact parameter for each
integer input. Its previous simultaneous-witness proof separated contact
lists by translations chosen after the inputs were known. We give fixed
translations that work for every integer tuple: distinct cosets modulo
the two-adic integers separate all the lists because their original
positions are two-adic integers. This reduces twelve contact coordinates
to one. A common square then removes all twelve private sum invariants,
at the cost of one coordinate and one equation. All substitutions are
affine, so no degree is increased. The resulting complete table has 786
variables and 685 quadratic residuals after the independently compiled
ceiling-power reduction; applying only the present transformation to the
preceding 837-variable Pell table gives 815 variables and 714 residuals.
The arithmetic premises, ordinary rational abc, ineffective sufficient
shift, and unresolved emended-source universality assertion are retained.
\end{abstract}

\section{A fixed-place separation lemma}

Normalize the rational valuation by \(v_p(p)=1\).

\begin{lemma}\label{lem:separation}
Let \(\mathcal A_0,\ldots,\mathcal A_{k-1}\) be finite sets in
\(\Q\cap\Z_p\). Choose \(e\ge0\) with \(p^e\ge k\), and set
\[
                  \delta_i=i/p^e\quad(0\le i<k).
\]
The sets \(\mathcal A_i-\delta_i\) are pairwise disjoint.
These same fixed translations work for every possible choice of the
sets in \(\Q\cap\Z_p\).
\end{lemma}
\begin{proof}
For \(i\ne j\), the nonzero integer \(i-j\) has absolute value less
than \(p^e\), so \(v_p(i-j)<e\). Thus
\(v_p(\delta_i-\delta_j)<0\), whereas every difference between
elements of \(\mathcal A_i\) and \(\mathcal A_j\) is in \(\Z_p\).
Equality of two translated positions would equate these quantities,
which is impossible.
\end{proof}

The argument also applies when every \(\mathcal A_i\) lies in
\(p^{-L}\Z_p\), for one fixed integer \(L\ge0\): use
\(\delta_i=i/p^{L+e}\).
There is no bound on the real size or cardinality of the sets.
Internal coincidences in an individual set are not removed by the lemma;
they must be dealt with by the original construction.

\section{Application to the ordinary-abc source}

We use the complete 56-coordinate interface of~\cite{abc56} and the
shared-reference theorem of~\cite{reference}. In that interface the
input is \(X\), the reference curve is
\[
 E:y^2=x^3-(11-6\sqrt2)x,\qquad
 R_0=(1+\sqrt2,2\sqrt2),
\]
and the integer-witness trace for \(X=n\ne0\) is
\[
 T_n=n^2\operatorname{Tr}x([n]R_0).
\]
Its proved two-adic trace lemma gives \(T_n\in\Z_2\). With the
sufficiently large fixed odd source parameter \(B\), put
\[
 M_n=B+2(n^8+T_n^2),\qquad a_n=M_n^4+16.
\]
Consequently \(M_n\in\Z_2^\times\), \(a_n\in\Z_2\), and in fact
\(a_n\equiv1\pmod8\) is a nonzero two-adic square. The use of
two-adic integrality here is restricted to constructing integer
witnesses. It is not assumed for arbitrary rational zeros.

Let
\[
 \mathcal R=\{-8,-7,-4,-1,1,4,7,8\},\qquad
 \mathcal A_n=\{\pm(a_n+r):r\in\mathcal R\}.
\]
Every position belongs to \(\Z_2\). Also \(a_n>16\), so these sixteen
positions are internally distinct. For \(k\) input blocks indexed by
\(0,\ldots,k-1\), replace their separate physical contact variables by
\begin{equation}\label{eq:affine}
                   z_i=z+\delta_i,
 \qquad\delta_i=i/2^e,\quad 2^e\ge k.
\end{equation}
The actual positions as functions of \(z\) are
\(\mathcal A_{n_i}-\delta_i\). Lemma~\ref{lem:separation} makes the
whole collection of \(16k\) positions distinct, uniformly in the
nonzero integer tuple.

\begin{theorem}\label{thm:tuple}
Fix a defining odd parameter \(B\) for the 56-coordinate interface,
under ordinary rational abc and its stated arithmetic premises.
In \(k\) copies identify the six reference coordinates exactly as
in~\cite{reference}, make substitutions~\eqref{eq:affine}, and use the
common contact square described below.
Every rational solution has integral inputs. Every tuple in
\((\Z\setminus\{0\})^k\) has a rational solution.
The resulting system has
\[
                   48k+8\quad\text{coordinates},\qquad
                   41k+5\quad\text{quadratic residuals},
\]
including its \(k\) distinguished inputs. Its sum of squares has degree
exactly four. The sufficient parameter threshold is unchanged.
\end{theorem}

\begin{proof}
For soundness, recover the old \(z_i\) from~\eqref{eq:affine}. Every
block is a solution of its unmodified 56-coordinate source. Its
integer-defining theorem therefore applies directly. In particular,
soundness requires neither two-adic integrality of the new center values
nor disjointness of their contact lists. No independent zero-input
simultaneous-completeness assertion is used.

For completeness fix nonzero integers \(n_0,\ldots,n_{k-1}\). Follow
the common-reference proof of~\cite{reference}, choosing
\[
 U_i=[n_i]H,\quad V_i=[n_i](H-R_0),\quad U_i-V_i=[n_i]R_0.
\]
The traces \(T_{n_i}\), centers \(a_{n_i}\), and all contact positions
are now fixed independently of the later reference choice. Choose a
common reference neighborhood avoiding the finitely many forbidden
points of every block, as in that proof. The finite tests are bounded
there. The fixed dyadic translations now separate the contact lists;
there is no input-dependent choice of offsets.

Enlarge finite avoidance to contain the supports of the input integers,
contact positions, all their nonzero differences, centers, chart
denominators, and the translations. This is a finite set fixed before
prime selection; it is not a change to the fixed exceptional set in
the source soundness theorem. Choose an integer \(D_0>0\) such that
\(D_0^2(d-\delta_i)\in\Z\) for every \(d\in\mathcal A_{n_i}\).
The retained finite-family prime supplier applies to
\[
 N-D_0^2(d-\delta_i)L^2,
 \qquad z=\frac{N}{D_0^2L^2}.
\]
These are distinct fixed rational contact positions, exactly the type
of family allowed in the previous shared-reference construction.
Use its same primary and opposite reference congruences, compatible
annihilation refinements, and real recurrence. At a selected prime
precisely one factor of precisely one contact polynomial vanishes
simply. The unchanged derivative identity gives contact-quotient
valuation one; the unchanged point preparation gives positive
valuation of the relevant detector error. Thus the selected native
pole valuations are the same as before.

For the other primes the source factorization and parity argument
applies separately to each block. Fixed translation denominators
introduce no missing case: they were explicitly included in finite
preparation. At any fixed finite place, take the denominator of \(z\)
sufficiently deep that
\[
                   z_i/z=1+\delta_i/z\longrightarrow1
\]
for all \(i\). The same limit holds as \(z\) tends to real infinity.
The source estimates consequently hold simultaneously:
\[
 u_i\sim z_i^2/a_{n_i},\qquad c_i\sim a_{n_i}/z_i^2,
 \qquad W_i-1=O(c_i),\qquad D_i=O(c_i^2).
\]
At two the centers are squares, so every actual pole input remains
in \(7\Q_2^{\times2}\), with arbitrarily large positive valuation.
At infinity each lies in \((0,1)\). Its reduced denominator is a
square by the same selected-prime and unselected-prime argument as
before. Apply the retained native-pole realization independently to
the actual inputs. This fills all blocks with their common reference
and with the exact affine contact parameters~\eqref{eq:affine}.

The shared-reference system has \(50k+6\) coordinates and
\(42k+4\) rows. Equation~\eqref{eq:affine} replaces its \(k\) contact
variables by one and is implemented by substitution, hence removes
\(k-1\) coordinates without adding equations. The common-square
substitution of the next section removes another \(k-1\) coordinates
and \(k-1\) equations, by an exact algebraic equivalence.
Affine substitution and constant denominator clearing preserve degree
at most two. An
unchanged input-square residual has nonzero quadratic part, so the
sum of squares has degree exactly four over \(\Q\).
\end{proof}

\section{Literal coordinate transformation and complete counts}

The source contact variable is local coordinate 48; it occurs only in
the row named \texttt{contact.difference}. If coordinates 49 and 50
are denoted by \(T_i,S_i\), that residual is
\[
                         S_i-T_i-2z_i^2.
\]
After~\eqref{eq:affine}, introduce one common coordinate \(q\), impose
\(q=z^2\), and make the affine substitutions
\begin{equation}\label{eq:sum-elimination}
               S_i=T_i+2q+4\delta_i z+2\delta_i^2.
\end{equation}
Delete all \(k\) private difference rows. Their substituted residuals
are identically \(2(q-z^2)\), so the one new square row is equivalent
to all of them. Every occurrence of \(S_i\) in the remaining source
rows is replaced by~\eqref{eq:sum-elimination}, including its squared
occurrences. Since this substitution is affine in the retained
coordinates, every row stays quadratic. Conversely a new solution
uniquely recovers the old \(S_i,z_i\), and an old solution satisfying
\eqref{eq:affine} has the unique extension \(q=z^2\).
This proves an exact equivalence on the whole rational zero set,
including \(z=0\), \(T_i=0\), and every other allowed boundary.

The eliminated local sum coordinate is 50 and the retained difference
coordinate is 49. The common square is named
\texttt{abc\_common\_contact\_square}. Each of the five affected row
types is transformed: \texttt{contact.sum}, \texttt{contact.A1},
\texttt{contact.A2}, \texttt{contact.partialCombined}, and
\texttt{contact.derivativeQuotient}. All other rows are copied
literally after renumbering. Nonzero constant denominators are cleared
and each changed row is divided by its integer content. For twelve
blocks we take \(e=4\), so all offsets have denominator dividing 16.
No variable
denominator or new guard is introduced.

For the emended Jones source the twelve certified input names are,
in the order assigned offsets,
\[
                b,g,t,m,w,s,h,c,\gamma,o,j,i.
\]
The last block's stored prefix remains \texttt{int\_i\_reciprocal},
but its distinguished ordinary input is \(i\), with the retained
equation \(i\,i_{\rm inverse}=1\). No reciprocal chart is changed.

Every positive integer source solution supplies nonzero values of all
twelve inputs. Theorem~\ref{thm:tuple} therefore gives completeness of
the source transfer. Conversely every new rational solution extends to
the previous complete table and hence to its integer source under the
same retained arithmetic premises. Thus the complete outer-instance
projection is unchanged, even though the present substitution is not
a bijection on every tuple of old auxiliary witnesses.

\begin{center}
\begin{tabular}{lrrr}
\toprule
Predecessor & Old coordinates & New coordinates & Rows\\
\midrule
Shared reference and Pell reduction &837&815&714\\
Further ceiling-power reduction &808&786&685\\
\bottomrule
\end{tabular}
\end{center}

The output is the sum of the squares of the complete listed integral
quadratic residuals. The integer instance occurs in the coefficients,
not among the displayed rational unknowns. Universality of the exact
emended source remains the separate assertion \(U_{\rm em}\), and
has not been established by the present theorem. Ordinary rational
abc, the source elliptic, height, prime-supplier and recurrence
premises, and the sufficiently large ineffective odd parameter are
all retained. The stored small \(B\) member is only an algebraic
fixture. No unconditional undecidability bound or global minimum is
asserted.

\section{Reproducibility}

The complete producer and independent checker are
\begin{center}\small
\nolinkurl{undecidability/joint_detector_v93/build_abc_fixed_offsets.py}\\
\nolinkurl{undecidability/joint_detector_v93/verify_abc_fixed_offsets.py}.
\end{center}
Both accept an explicit predecessor path. The producer uses rational
coefficient dictionaries; the checker independently uses symbolic
affine substitution and primitive polynomial normalization. The
checker reconstructs every changed row, compares every unchanged row,
checks complete occurrence lists, exact counts, coefficient domains,
parameter qualifications and file hashes. The general separation and
arithmetic-completeness arguments are proofs above; finite valuation
examples in the checker are only illustrations.

\begin{thebibliography}{9}
\bibitem{abc56} \(\Z\Q\) Diophantine Project,
\emph{Paired contact coordinates: an ordinary-abc quartic in 56 variables},
19 November 2026,
\nolinkurl{papers/quartic_existence/paired_contact56_20261119_v81.tex}.
\bibitem{reference} \(\Z\Q\) Diophantine Project,
\emph{A shared-reference quartic checkpoint with an explicit
source-universality obligation}, 23 November 2026,
\nolinkurl{undecidability/papers/shared_quartic_checkpoint.tex}.
\end{thebibliography}
\end{document}
